\documentclass[10pt,a4paper]{amsart}
\usepackage[margin=19mm]{geometry}
\usepackage{amsmath,amssymb,mathtools}
\usepackage[scr=rsfs,cal=euler]{mathalfa}
\usepackage{xfrac}
\usepackage[unicode,hidelinks,bookmarks=false]{hyperref}
\newcommand{\ie}{i.e.\ }
\newcommand{\cf}{cf.\ }
\newcommand{\resp}{respectively\ }
\newcommand{\Subsection}{\subsection}
\providecommand{\nobreakdash}{\nobreak\hbox{-}}
\setlength{\emergencystretch}{2em}
\multlinegap0pt
\usepackage{bm}
\usepackage{bbm}
\usepackage{dsfont}
\usepackage{xspace}
\usepackage{cancel}
\usepackage{mathtools}
\usepackage[shortlabels]{enumitem}
\usepackage{amsthm}
\makeatletter
\@ifundefined{theorem}{\newtheorem{theorem}{Theorem}}{}
\@ifundefined{definition}{\newtheorem{definition}[theorem]{Definition}}{}
\@ifundefined{lemma}{\newtheorem{lemma}[theorem]{Lemma}}{}
\makeatother
\title[Tessellations and Speiser graphs arising from meromorphic fu]{Tessellations and Speiser graphs arising from meromorphic functions on simply connected Riemann surfaces}
\author{Alvaro Alvarez-Parrilla, Jesus Muci\~no-Raymundo}
\date{Source: arXiv:2602.00416v1 (CC BY 4.0)}
\begin{document}
\maketitle

\noindent\emph{Source and licence.} Tessellations and Speiser graphs arising from meromorphic functions on simply connected Riemann surfaces. Version arXiv:2602.00416v1 (CC BY 4.0), distributed under the Creative Commons Attribution 4.0 International licence, as recorded in the arXiv licence metadata for this exact version and shown on the abstract page https://arxiv.org/abs/2602.00416v1. Full rights notice: licenses/arxiv-2602.00416-CC-BY-4.0.txt.\par\smallskip

\noindent\textit{Editorial scope.} Counted unit: the Lemma whose source label is \texttt{lem:A-mapa-admite-etiquetados} in the article (source0.tex lines 3524--3529, source label \texttt{lem:A-mapa-admite-etiquetados}), delivered together with the article's own complete original proof of it (source0.tex lines 3531--3576, one \texttt{proof} environment of 46 lines). No mathematics is written, repaired or re-derived: statement and proof are the article's own text line for line. Required context retained uncounted: def:de-teselacion (lines 2593--2652); definicion-de-A-map (lines 2735--2782); def:etiquetado-consistente (lines 3002--3031). Every definition, display and label of those lines is retained unchanged. Omitted from this package and not redistributed: 1--2592, 2653--2734, 2783--3001, 3032--3523, 3530--3530, 3577--9017. Nothing inside a retained excerpt is omitted or altered: every retained body is delivered exactly as the article's lines are written, so no omission receipt of any kind accompanies this package. Citations: the retained excerpts carry no citation command at all, so no bibliography entry of the article is transcribed and no reference is redistributed. Reference adaptation: none was needed and none was made; every reference inside the delivered unit resolves inside it, so no unresolved reference or citation is produced. Printed slips kept literal: no printed word, symbol, hypothesis or number of the counted statement or of its proof was altered. Layout and class: the article's own class is \texttt{elsarticle} with packages including amssymb, amsmath, mathtools, enumitem, mathrsfs, stackrel, graphicx, upgreek, hyperref, placeins, flafter, xcolor, ulem, tikz; the wrapper is the lane's pinned \texttt{amsart} editorial wrapper at 10pt on A4, which restates the theorem-like environments the retained text needs and adds the article's own macro definitions that the retained text needs (none), with extra wrapper packages (bm, bbm, dsfont, xspace, cancel, mathtools, enumitem). The delivered numbering is the unit's own. Assets: no retained span contains a figure environment, a graphics-inclusion command, a PDF-page inclusion, a table or a TikZ picture; no image file is copied into this package, the package declares no asset dependency and no image file is redistributed with it. Licence: version arXiv:2602.00416v1 of this article is distributed under the Creative Commons Attribution 4.0 International licence, as recorded in the arXiv licence metadata for this exact version and shown on the abstract page https://arxiv.org/abs/2602.00416v1. A full rights notice is delivered at licenses/arxiv-2602.00416-CC-BY-4.0.txt. Characters: every retained excerpt is pure ASCII, so the delivered engine, XeLaTeX with its default Latin Modern text font, reports no missing character.
\begin{definition}
\label{def:de-teselacion}
\
\begin{enumerate}[label=\arabic*),leftmargin=*] 
\item
A  \emph{tessellation} of a
surface $\Omega_z$ is a collection 
of alternating colored tiles
\begin{equation}
\label{teselacion}
\mathscr{T}
=
\underbrace{
T_1 \cup \ldots \cup T_\alpha \cup \ldots  }_{ n\text{ blue tiles} }
\cup 
\underbrace{ T^\prime_1 \cup  
\ldots  \cup T^{\prime}_\alpha  \cup \ldots }_{n \text{ grey tiles}} \subset 
\Omega_z \cup \partial_{\mathcal{I}} \Omega_z, 
\ \ \ 
2 \leq n \leq \infty,
\end{equation}

\noindent
where the \emph{tiles}
$\{ T_{\alpha}, \, T_{\alpha}^\prime \}$
are open Jordan domains, 
such that: 

\begin{enumerate}[label=\roman*),leftmargin=*]
\item
The union of their closures 
$\cup_{{\alpha}}
\big( \overline{T_{\alpha}} 
\cup
\overline{T_{\alpha}^\prime } \big)$
is $\Omega_z \cup \partial_{\mathcal{I}} \Omega_z$.

\item 
The boundary of the closure of each tile 
$\partial\overline{T_\alpha}$ (resp. 
$\partial\overline{ T_{\alpha}^\prime}$) has 
$\rho$ vertices and $\rho$ edges, where $2 \leq \rho \leq {\tt q}$
and $\rho$ depends on the particular tile. 
 


\item
If the intersection of the closures of any
two tiles is non--empty, 
then it consists of a finite number of 
simple paths (edges) and their extreme points (vertices).
\end{enumerate}

\item
If all the tiles $\{ T_{\alpha} \}$
(resp. $\{T_{\alpha}^\prime \}$) have 
the same number of vertices and edges,
we shall say that the tessellation is \emph{homogeneous}.
\end{enumerate}
\end{definition}

\begin{definition}
\label{definicion-de-A-map}
\begin{enumerate}[label=\arabic*),leftmargin=*] 
\item
An {\it ${\tt A}$--map} $\widehat{\Gamma}_{\tt q}$ is an 
oriented, connected graph 
embedded in  $\Omega_z \cup \partial_{\mathcal{I}} \Omega_z$, 
with vertices $V(\widehat{\Gamma}_{\tt q})$ of
infinite or 
even valence  greater than or equal to $2$
and edges $E (\widehat{\Gamma}_{\tt q})$, such that:

\begin{enumerate}[label=\roman*),leftmargin=*]
\item
The subset of vertices of valence greater
than or equal to 4 is non empty. 

\item
If we forget all the vertices
of valence 2 of $\widehat{\Gamma}_{\tt q}$, then
we obtain a ${\tt t}$--graph $\Gamma$ such that:
%
\begin{equation}
\label{eq:tessellation}
%
\mathscr{T}(\widehat{\Gamma}_{\tt q}) \doteq %&
\mathscr{T}(\Gamma)
\end{equation}



\noindent
is, set theoretically, a tessellation as in Definition 
\ref{def:de-teselacion}.
\end{enumerate}

\item 
The boundary 
$\partial \overline{T_\alpha}$ 
(resp. $\partial \overline{T_{\alpha}^\prime}$) of  
each tile consists of exactly $\tt q$ vertices and $\tt q$ edges 
of $\widehat{\Gamma}_{\tt q}$,
{\it i.e.} the tessellation $\mathscr{T}(\widehat{\Gamma}_{\tt q})$ is homogeneous.

\item 
We shall say that $\Gamma$ and $\widehat{\Gamma}_{\tt q}$, as in (ii) above, are \emph{compatible}.
\end{enumerate}
\end{definition}

\begin{definition}
\label{def:etiquetado-consistente}
A {\it consistent $q$--labelling 

\centerline{$\mathcal{L}_{\mathcal{W}_{\tt q}}: V(\Gamma) \longrightarrow \mathcal{W}_{\tt q}, 
\ \ \
{\tt q} \geq 2$, } 

\noindent 
for a ${\tt t}$--graph $\Gamma$} satisfies the following conditions:

\begin{enumerate}[label=\roman*),leftmargin=*]
\item
For each blue tile $T_\alpha$ of the tessellation
$\mathscr{T}(\Gamma)$, 
if $\{ z_\iota \}$ are the vertices of its
boundary $\partial \overline{T_\alpha}$, ordered 
with cyclic anti--clockwise sense, then the 
labels (values) $\{ \mathcal{L}_{\mathcal{W}_{\tt q}} ( z_\iota ) \} \subset \mathcal{W}_{\tt q}$
appear exactly once and with the same cyclic 
order provided by $\mathcal{L}_\gamma$.

\item
Each label (value) ${\tt w}_{j} \in \mathcal{W}_{\tt q}$ appears under
$\mathcal{L}_{\mathcal{W}_{\tt q}}$ for at least one
vertex $z_\iota \in V(\Gamma)$ of
$\Gamma$, which by definition have
valence greater than or equal to $4$.
\end{enumerate}  
\end{definition}

\begin{lemma}
\label{lem:A-mapa-admite-etiquetados}
An $\tt A$--map $\widehat{\Gamma}_{\tt q}$ supports at least one consistent 
${\tt q}_{\rm o}$--labelling $\mathcal{L}_{\mathcal{W}_{{\tt q}_{\rm o}}}$, for 
$2 \leq {\tt q}_{\text{min}} \leq {\tt q}_{\rm o} \leq {\tt q} \leq {\tt q}_{\text{max}}$.
\end{lemma}

\begin{proof}
Given the $\tt A$--map $\widehat{\Gamma}_{\tt q}$, 
let $\Gamma$ be the $\tt t$--graph, alluded to in Definition \ref{definicion-de-A-map}.ii
(in other words $\Gamma$ and $\widehat{\Gamma}_{\tt q}$ are compatible).
Note that ${\tt q}_{\text{min}} \leq {\tt q} \leq {\tt q}_{\text{max}}$ as in the above Remark.
The existence of a consistent ${\tt q}$--labelling is as follows:
\begin{enumerate}[label=\arabic*),leftmargin=*]
\item Choose ${\tt q}_{\rm o} = {\tt q}$,

\item 
assign the labels $\mathcal{W}_{{\tt q}_{\rm o}}$ to the ${\tt q}_{\rm o}$ vertices of any 
blue tile $T_\alpha$ of $\widehat{\Gamma}_{\tt q}$, in an anticlockwise order,

\item
propagate the labelling to all the neighbor grey tiles, 
this can be done since
the tessellation $\mathcal{T}(\widehat{\Gamma}_{\tt q})$ is homogeneous, \emph{i.e.}\, 
its tiles are topological ${\tt q}_{\rm o}$--gons,

\item
continue as above to all the tiles of $\mathcal{T}(\widehat{\Gamma}_{{\tt q}_{\rm o}})$,
using the fact that $\Omega_z$ is simply connected.
\end{enumerate}
This provides a labelling 

\centerline{$
\mathcal{L}_{\mathcal{W}_{{\tt q}_{\rm o}}}: V(\widehat{\Gamma}_{{\tt q}_{\rm o}}) \longrightarrow \mathcal{W}_{{\tt q}_{\rm o}}, 
$}

\noindent
to $\widehat{\Gamma}_{{\tt q}_{\rm o}}$.
Clearly, $\mathcal{L}_{\mathcal{W}_{{\tt q}_{\rm o}}}$ satisfies condition (i) of a consistent ${\tt q}_{\rm o}$--labelling,
see Definition \ref{def:etiquetado-consistente}. 

\noindent
If the choice of ${\tt q}_{\rm o}$ satisfies condition (ii) then we are done.

\noindent
Otherwise, there is a label (value), say ${\tt w}_{{\tt j}_{\rm o}}$, that does not appear under 
$\mathcal{L}_{\mathcal{W}_{{\tt q}_{\rm o}}}$ for a vertex of $\widehat{\Gamma}_{{\tt q}_{\rm o}}$
of valence greater than or equal to 4. 
We can erase this label from the cyclic order $\mathcal{W}_{{\tt q}_{\rm o}}$, 
obtaining a new cyclic order with ${\tt q}_{\rm o} - 1$ distinct values.
Now, go back to (2) with a new ${\tt q}_{\rm o}$ being one less than the previous one and repeat the process.
Since the original $\tt A$--map $\widehat{\Gamma}_{\tt q}$ has a non empty subset of vertices of valence greater than or equal to 4, this process eventually stops at a ${\tt q}_{\rm o} \geq {\tt q}_{\text{min}}$.
\end{proof}

\end{document}
